\section{Theoretical Framework}

Skeleton-based human activity recognition (HAR) describes a person as a
sequence of body keypoints rather than raw pixels. This reduces each frame to
a few dozen values and makes the input less sensitive to clothing, background
and lighting. The most accurate skeleton classifiers are deep networks. They
lead benchmarks such as NTU RGB+D, but they generally assume GPU hardware. A
lighter, well-established alternative computes hand-crafted statistics over a
sliding temporal window and classifies them with a random forest
\cite{breiman2001random}. This approach is fast, works with modest amounts of
data, keeps its features interpretable and is small enough to run on the device
itself.

Existing work highlights two uncertainties for this approach. First, windowed
activity data contains near-duplicate samples from the same person, so
record-wise train/test splits overestimate performance; subject-wise
validation is needed to approximate the real use case
\cite{saeb2017need}. Second, public datasets are recorded in specific settings,
so a model that scores well on one corpus may not transfer to footage from
another, such as a robot's camera in a living room. The innovative character
of this part of the work is therefore not the classifier itself. It is the
combination of training on published skeleton datasets, evaluating with subject
separation, exporting the forest to a NumPy-only format and connecting it
through the robot's SDK. As a result, recognition runs on the Reachy Mini's CPU
and directly drives expressive feedback.